\\newpage
\\section{问题分析}

\\subsection{问题一分析}
本题要求对智能电网负荷和新能源发电数据进行\\textbf{短期预测}。负荷数据具有明显的日周期性和周周期性，同时受气温、节假日等因素影响。传统ARIMA模型擅长捕捉线性趋势，但对非线性模式建模能力有限；LSTM可以处理长序列依赖，但训练成本高且难以捕捉长期全局依赖。

\\textbf{解题思路}：
\\begin{enumerate}
    \\item 采用STL分解方法将原始序列分解为趋势项、季节项和残差项
    \\item 趋势项和季节项采用Transformer模型的Self-Attention机制进行建模
    \\item 残差项采用ARIMA进行线性修正
    \\item 最终预测值 = Transformer预测值 + ARIMA残差修正值
\\end{enumerate}

\\subsection{问题二分析}
本题需要建立\\textbf{强化学习调度环境}。这是典型的序贯决策问题：每个时刻的调度决策会影响后续状态的演化。我们需要：
\\begin{itemize}
    \\item 合理定义状态空间，包含负荷、发电、储能等关键变量
    \\item 定义离散化的动作空间，表示各发电机组的出力调整
    \\item 设计综合奖励函数，平衡成本、消纳和满意度
\\end{itemize}
采用Q-Learning算法，通过试错学习最优调度策略。

\\subsection{问题三分析}
本题是多目标优化问题，三个目标之间存在\\textbf{冲突关系}：降低调度成本往往意味着减少备用容量，可能影响供电可靠性；提高新能源消纳率可能需要增加储能投入，推高成本。我们采用NSGA-II算法求解Pareto前沿，并通过熵权-TOPSIS方法进行方案排序和选择。
